\section{Intersecting with Chern classes}
\label{section-intersecting-chern-classes}

\noindent
In this section we define Chern classes of vector bundles on $X$ as
bivariant classes on $X$, see Lemma \ref{lemma-cap-cp-bivariant}
and the discussion following this lemma. Our construction follows the familiar
pattern of first defining the operation on prime cycles and then
summing. In Lemma \ref{lemma-determine-intersections} we show
that the result is determined by the usual formula on the associated
projective bundle. Next, we show that capping with Chern classes
passes through rational equivalence, commutes with proper pushforward,
commutes with flat pullback, and commutes with the Gysin maps for
inclusions of effective Cartier divisors. These lemmas could have been
avoided by directly using the characterization in
Lemma \ref{lemma-determine-intersections} and using
Lemma \ref{lemma-push-proper-bivariant}; the reader who wishes to
see this worked out should consult
Chow Groups of Spaces, Lemma \ref{spaces-chow-lemma-segre-classes}.

\begin{definition}
\label{definition-cap-chern-classes}
Let $(S, \delta)$ be as in Situation \ref{situation-setup}.
Let $X$ be locally of finite type over $S$.
Let $\mathcal{E}$ be a finite locally free sheaf of rank $r$ on $X$.
We define, for every integer $k$ and any $0 \leq j \leq r$,
an operation
$$
c_j(\mathcal{E}) \cap - : Z_k(X) \to \CH_{k - j}(X)
$$
called {\it intersection with the $j$th Chern class of $\mathcal{E}$}.
\begin{enumerate}
\item Given an integral closed subscheme $i : W \to X$ of $\delta$-dimension
$k$ we define
$$
c_j(\mathcal{E}) \cap [W] = i_*(c_j({i^*\mathcal{E}}) \cap [W])
\in
\CH_{k - j}(X)
$$
where $c_j({i^*\mathcal{E}}) \cap [W]$ is as defined in
Definition \ref{definition-chern-classes}.
\item For a general $k$-cycle $\alpha = \sum n_i [W_i]$ we set
$$
c_j(\mathcal{E}) \cap \alpha = \sum n_i c_j(\mathcal{E}) \cap [W_i]
$$
\end{enumerate}
\end{definition}

\noindent
If $\mathcal{E}$ has rank $1$ then this agrees with our
previous definition (Definition \ref{definition-cap-c1})
by Lemma \ref{lemma-first-chern-class}.

\begin{lemma}
\label{lemma-determine-intersections}
Let $(S, \delta)$ be as in Situation \ref{situation-setup}.
Let $X$ be locally of finite type over $S$.
Let $\mathcal{E}$ be a finite locally free sheaf of rank $r$ on $X$.
Let $(\pi : P \to X, \mathcal{O}_P(1))$ be the projective bundle
associated to $\mathcal{E}$.
For $\alpha \in Z_k(X)$ the elements
$c_j(\mathcal{E}) \cap \alpha$ are the unique elements
$\alpha_j$ of $\CH_{k - j}(X)$
such that $\alpha_0 = \alpha$ and
$$
\sum\nolimits_{i = 0}^r
(-1)^i c_1(\mathcal{O}_P(1))^i \cap
\pi^*(\alpha_{r - i}) = 0
$$
holds in the Chow group of $P$.
\end{lemma}

\begin{proof}
The uniqueness of $\alpha_0, \ldots, \alpha_r$ such that
$\alpha_0 = \alpha$ and such that
the displayed equation holds follows from
the projective space bundle formula
Lemma \ref{lemma-chow-ring-projective-bundle}.
The identity holds by definition for $\alpha = [W]$ where $W$
is an integral closed subscheme of $X$.
For a general $k$-cycle $\alpha$ on $X$ write
$\alpha = \sum n_a[W_a]$ with $n_a \not = 0$, and
$i_a : W_a \to X$ pairwise distinct integral closed subschemes.
Then the family $\{W_a\}$ is locally finite on $X$.
Set $P_a = \pi^{-1}(W_a) = \mathbf{P}(\mathcal{E}|_{W_a})$.
Denote $i'_a : P_a \to P$ the corresponding closed immersions.
Consider the fibre product diagram
$$
\xymatrix{
P' \ar@{=}[r] \ar[d]_{\pi'} &
\coprod P_a \ar[d]_{\coprod \pi_a} \ar[r]_{\coprod i'_a} &
P \ar[d]^\pi \\
X' \ar@{=}[r] &
\coprod W_a \ar[r]^{\coprod i_a} &
X
}
$$
The morphism $p : X' \to X$ is proper. Moreover
$\pi' : P' \to X'$ together with the invertible sheaf
$\mathcal{O}_{P'}(1) = \coprod \mathcal{O}_{P_a}(1)$
which is also the pullback of $\mathcal{O}_P(1)$
is the projective bundle associated to
$\mathcal{E}' = p^*\mathcal{E}$. By definition
$$
c_j(\mathcal{E}) \cap [\alpha]
=
\sum i_{a, *}(c_j(\mathcal{E}|_{W_a}) \cap [W_a]).
$$
Write $\beta_{a, j} = c_j(\mathcal{E}|_{W_a}) \cap [W_a]$
which is an element of $\CH_{k - j}(W_a)$. We have
$$
\sum\nolimits_{i = 0}^r
(-1)^i c_1(\mathcal{O}_{P_a}(1))^i \cap \pi_a^*(\beta_{a, r - i}) = 0
$$
for each $a$ by definition. Thus clearly we have
$$
\sum\nolimits_{i = 0}^r
(-1)^i c_1(\mathcal{O}_{P'}(1))^i \cap (\pi')^*(\beta_{r - i}) = 0
$$
with $\beta_j = \sum n_a\beta_{a, j} \in \CH_{k - j}(X')$. Denote
$p' : P' \to P$ the morphism $\coprod i'_a$.
We have $\pi^*p_*\beta_j = p'_*(\pi')^*\beta_j$
by Lemma \ref{lemma-flat-pullback-proper-pushforward}.
By the projection formula of Lemma \ref{lemma-pushforward-cap-c1}
we conclude that
$$
\sum\nolimits_{i = 0}^r
(-1)^i c_1(\mathcal{O}_P(1))^i \cap \pi^*(p_*\beta_j) = 0
$$
Since $p_*\beta_j$ is a representative of $c_j(\mathcal{E}) \cap \alpha$
we win.
\end{proof}

\noindent
We will consistently use this characterization of Chern classes
to prove many more properties.

\begin{lemma}
\label{lemma-cap-chern-class-factors-rational-equivalence}
Let $(S, \delta)$ be as in Situation \ref{situation-setup}.
Let $X$ be locally of finite type over $S$.
Let $\mathcal{E}$ be a finite locally free sheaf of rank $r$ on $X$.
If $\alpha \sim_{rat} \beta$ are rationally equivalent $k$-cycles
on $X$ then $c_j(\mathcal{E}) \cap \alpha = c_j(\mathcal{E}) \cap \beta$
in $\CH_{k - j}(X)$.
\end{lemma}

\begin{proof}
By Lemma \ref{lemma-determine-intersections} the elements
$\alpha_j = c_j(\mathcal{E}) \cap \alpha$, $j \geq 1$ and
$\beta_j = c_j(\mathcal{E}) \cap \beta$, $j \geq 1$ are uniquely determined
by the {\it same} equation in the chow group of the projective
bundle associated to $\mathcal{E}$. (This of course relies on the fact that
flat pullback is compatible with rational equivalence, see
Lemma \ref{lemma-flat-pullback-rational-equivalence}.) Hence they are equal.
\end{proof}

\noindent
In other words capping with Chern classes of
finite locally free sheaves factors through rational equivalence
to give maps
$$
c_j(\mathcal{E}) \cap - : \CH_k(X) \to \CH_{k - j}(X).
$$
Our next task is to show that Chern classes are bivariant classes, see
Definition \ref{definition-bivariant-class}.

\begin{lemma}
\label{lemma-pushforward-cap-cj}
Let $(S, \delta)$ be as in Situation \ref{situation-setup}.
Let $X$, $Y$ be locally of finite type over $S$.
Let $\mathcal{E}$ be a finite locally free sheaf of rank $r$ on $X$.
Let $p : X \to Y$ be a proper morphism.
Let $\alpha$ be a $k$-cycle on $X$.
Let $\mathcal{E}$ be a finite locally free sheaf on $Y$.
Then
$$
p_*(c_j(p^*\mathcal{E}) \cap \alpha) = c_j(\mathcal{E}) \cap p_*\alpha
$$
\end{lemma}

\begin{proof}
Let $(\pi : P \to Y, \mathcal{O}_P(1))$ be the projective bundle associated
to $\mathcal{E}$. Then $P_X = X \times_Y P$ is the projective bundle associated
to $p^*\mathcal{E}$ and $\mathcal{O}_{P_X}(1)$ is the pullback of
$\mathcal{O}_P(1)$. Write $\alpha_j = c_j(p^*\mathcal{E}) \cap \alpha$, so
$\alpha_0 = \alpha$. By Lemma \ref{lemma-determine-intersections} we have
$$
\sum\nolimits_{i = 0}^r
(-1)^i c_1(\mathcal{O}_P(1))^i \cap
\pi_X^*(\alpha_{r - i}) = 0
$$
in the chow group of $P_X$. Consider the fibre product diagram
$$
\xymatrix{
P_X \ar[r]_-{p'} \ar[d]_{\pi_X} & P \ar[d]^\pi \\
X \ar[r]^p & Y
}
$$
Apply proper pushforward $p'_*$
(Lemma \ref{lemma-proper-pushforward-rational-equivalence})
to the displayed equality above. Using
Lemmas \ref{lemma-pushforward-cap-c1} and
\ref{lemma-flat-pullback-proper-pushforward} we obtain
$$
\sum\nolimits_{i = 0}^r
(-1)^i c_1(\mathcal{O}_P(1))^i \cap
\pi^*(p_*\alpha_{r - i}) = 0
$$
in the chow group of $P$. By the characterization of
Lemma \ref{lemma-determine-intersections} we conclude.
\end{proof}

\begin{lemma}
\label{lemma-flat-pullback-cap-cj}
Let $(S, \delta)$ be as in Situation \ref{situation-setup}.
Let $X$, $Y$ be locally of finite type over $S$.
Let $\mathcal{E}$ be a finite locally free sheaf of rank $r$ on $Y$.
Let $f : X \to Y$ be a flat morphism of relative dimension $r$.
Let $\alpha$ be a $k$-cycle on $Y$.
Then
$$
f^*(c_j(\mathcal{E}) \cap \alpha) = c_j(f^*\mathcal{E}) \cap f^*\alpha
$$
\end{lemma}

\begin{proof}
Write $\alpha_j = c_j(\mathcal{E}) \cap \alpha$, so $\alpha_0 = \alpha$.
By Lemma \ref{lemma-determine-intersections} we have
$$
\sum\nolimits_{i = 0}^r
(-1)^i c_1(\mathcal{O}_P(1))^i \cap
\pi^*(\alpha_{r - i}) = 0
$$
in the chow group of the projective bundle
$(\pi : P \to Y, \mathcal{O}_P(1))$
associated to $\mathcal{E}$. Consider the fibre product diagram
$$
\xymatrix{
P_X = \mathbf{P}(f^*\mathcal{E}) \ar[r]_-{f'} \ar[d]_{\pi_X} &
P \ar[d]^\pi \\
X \ar[r]^f & Y
}
$$
Note that $\mathcal{O}_{P_X}(1)$ is the pullback of $\mathcal{O}_P(1)$.
Apply flat pullback $(f')^*$
(Lemma \ref{lemma-flat-pullback-rational-equivalence}) to the displayed
equation above. By Lemmas \ref{lemma-flat-pullback-cap-c1} and
\ref{lemma-compose-flat-pullback} we see that
$$
\sum\nolimits_{i = 0}^r
(-1)^i c_1(\mathcal{O}_{P_X}(1))^i \cap
\pi_X^*(f^*\alpha_{r - i}) = 0
$$
holds in the chow group of $P_X$. By the characterization of
Lemma \ref{lemma-determine-intersections} we conclude.
\end{proof}

\begin{lemma}
\label{lemma-cap-chern-class-commutes-with-gysin}
Let $(S, \delta)$ be as in Situation \ref{situation-setup}.
Let $X$ be locally of finite type over $S$.
Let $\mathcal{E}$ be a finite locally free sheaf of rank $r$ on $X$.
Let $(\mathcal{L}, s, i : D \to X)$ be as in
Definition \ref{definition-gysin-homomorphism}.
Then $c_j(\mathcal{E}|_D) \cap i^*\alpha = i^*(c_j(\mathcal{E}) \cap \alpha)$
for all $\alpha \in \CH_k(X)$.
\end{lemma}

\begin{proof}
Write $\alpha_j = c_j(\mathcal{E}) \cap \alpha$, so $\alpha_0 = \alpha$.
By Lemma \ref{lemma-determine-intersections} we have
$$
\sum\nolimits_{i = 0}^r
(-1)^i c_1(\mathcal{O}_P(1))^i \cap
\pi^*(\alpha_{r - i}) = 0
$$
in the chow group of the projective bundle
$(\pi : P \to X, \mathcal{O}_P(1))$
associated to $\mathcal{E}$. Consider the fibre product diagram
$$
\xymatrix{
P_D = \mathbf{P}(\mathcal{E}|_D) \ar[r]_-{i'} \ar[d]_{\pi_D} &
P \ar[d]^\pi \\
D \ar[r]^i & X
}
$$
Note that $\mathcal{O}_{P_D}(1)$ is the pullback of $\mathcal{O}_P(1)$.
Apply the Gysin map $(i')^*$ (Lemma \ref{lemma-gysin-factors}) to the
displayed equation above.
Applying Lemmas \ref{lemma-gysin-commutes-cap-c1} and
\ref{lemma-gysin-flat-pullback} we obtain
$$
\sum\nolimits_{i = 0}^r
(-1)^i c_1(\mathcal{O}_{P_D}(1))^i \cap
\pi_D^*(i^*\alpha_{r - i}) = 0
$$
in the chow group of $P_D$.
By the characterization of Lemma \ref{lemma-determine-intersections}
we conclude.
\end{proof}

\noindent
At this point we have enough material to be able to prove that
capping with Chern classes defines a bivariant class.

\begin{lemma}
\label{lemma-cap-cp-bivariant}
Let $(S, \delta)$ be as in Situation \ref{situation-setup}.
Let $X$ be locally of finite type over $S$.
Let $\mathcal{E}$ be a locally free $\mathcal{O}_X$-module
of rank $r$. Let $0 \leq p \leq r$.
Then the rule that to $f : X' \to X$ assigns
$c_p(f^*\mathcal{E}) \cap - : \CH_k(X') \to \CH_{k - p}(X')$
is a bivariant class of degree $p$.
\end{lemma}

\begin{proof}
Immediate from Lemmas
\ref{lemma-cap-chern-class-factors-rational-equivalence},
\ref{lemma-pushforward-cap-cj},
\ref{lemma-flat-pullback-cap-cj}, and
\ref{lemma-cap-chern-class-commutes-with-gysin}
and Definition \ref{definition-bivariant-class}.
\end{proof}

\noindent
This lemma allows us to define the Chern classes of a finite
locally free module as follows.

\begin{definition}
\label{definition-chern-classes-final}
Let $(S, \delta)$ be as in Situation \ref{situation-setup}.
Let $X$ be locally of finite type over $S$.
Let $\mathcal{E}$ be a locally free $\mathcal{O}_X$-module
of rank $r$. For $i = 0, \ldots, r$ the {\it $i$th Chern class}
of $\mathcal{E}$ is the bivariant class
$c_i(\mathcal{E}) \in A^i(X)$ of degree $i$
constructed in Lemma \ref{lemma-cap-cp-bivariant}. The
{\it total Chern class} of $\mathcal{E}$ is the formal sum
$$
c(\mathcal{E}) = 
c_0(\mathcal{E}) + c_1(\mathcal{E}) + \ldots + c_r(\mathcal{E})
$$
which is viewed as a nonhomogeneous bivariant class on $X$.
\end{definition}

\noindent
By the remark following Definition \ref{definition-cap-chern-classes}
if $\mathcal{E}$ is invertible, then this definition agrees with
Definition \ref{definition-first-chern-class}.
Next we see that Chern classes are in the center of the bivariant
Chow cohomology ring $A^*(X)$.

\begin{lemma}
\label{lemma-cap-commutative-chern}
Let $(S, \delta)$ be as in Situation \ref{situation-setup}.
Let $X$ be locally of finite type over $S$.
Let $\mathcal{E}$ be a locally free $\mathcal{O}_X$-module of rank $r$.
Then
\begin{enumerate}
\item $c_j(\mathcal{E}) \in A^j(X)$ is in the center of $A^*(X)$ and
\item if $f : X' \to X$ is locally of finite type and $c \in A^*(X' \to X)$,
then $c \circ c_j(\mathcal{E}) = c_j(f^*\mathcal{E}) \circ c$.
\end{enumerate}
In particular, if $\mathcal{F}$ is a second locally free
$\mathcal{O}_X$-module on $X$ of rank $s$, then
$$
c_i(\mathcal{E}) \cap c_j(\mathcal{F}) \cap \alpha
=
c_j(\mathcal{F}) \cap c_i(\mathcal{E}) \cap \alpha
$$
as elements of $\CH_{k - i - j}(X)$ for all $\alpha \in \CH_k(X)$.
\end{lemma}

\begin{proof}
It is immediate that (2) implies (1).
Let $\alpha \in \CH_k(X)$. Write $\alpha_j = c_j(\mathcal{E}) \cap \alpha$, so
$\alpha_0 = \alpha$. By Lemma \ref{lemma-determine-intersections} we have
$$
\sum\nolimits_{i = 0}^r
(-1)^i c_1(\mathcal{O}_P(1))^i \cap
\pi^*(\alpha_{r - i}) = 0
$$
in the chow group of the projective bundle
$(\pi : P \to Y, \mathcal{O}_P(1))$
associated to $\mathcal{E}$. Denote $\pi' : P' \to X'$ the base change
of $\pi$ by $f$. Using Lemma \ref{lemma-c1-center} and
the properties of bivariant classes we obtain
\begin{align*}
0 & = c \cap \left(\sum\nolimits_{i = 0}^r
(-1)^i c_1(\mathcal{O}_P(1))^i \cap
\pi^*(\alpha_{r - i})\right) \\
& =
\sum\nolimits_{i = 0}^r
(-1)^i c_1(\mathcal{O}_{P'}(1))^i \cap
(\pi')^*(c \cap \alpha_{r - i})
\end{align*}
in the Chow group of $P'$ (calculation omitted).
Hence we see that $c \cap \alpha_j$ is
equal to $c_j(f^*\mathcal{E}) \cap (c \cap \alpha)$ by the characterization
of Lemma \ref{lemma-determine-intersections}.
This proves the lemma.
\end{proof}

\begin{remark}
\label{remark-extend-to-finite-locally-free}
Let $(S, \delta)$ be as in Situation \ref{situation-setup}.
Let $X$ be locally of finite type over $S$.
Let $\mathcal{E}$ be a finite locally free $\mathcal{O}_X$-module.
If the rank of $\mathcal{E}$ is not constant then we can
still define the Chern classes of $\mathcal{E}$. Namely, in this
case we can write
$$
X = X_0 \amalg X_1 \amalg X_2 \amalg \ldots
$$
where $X_r \subset X$ is the open and closed subspace where
the rank of $\mathcal{E}$ is $r$. By 
Lemma \ref{lemma-disjoint-decomposition-bivariant}
we have $A^p(X) = \prod A^p(X_r)$.
Hence we can define $c_p(\mathcal{E})$ to be the
product of the classes $c_p(\mathcal{E}|_{X_r})$ in $A^p(X_r)$.
Explicitly, if $X' \to X$ is a morphism locally of finite type,
then we obtain by pullback a corresponding decomposition of $X'$
and we find that
$$
\CH_*(X') = \prod\nolimits_{r \geq 0} \CH_*(X'_r)
$$
by our definitions. Then $c_p(\mathcal{E}) \in A^p(X)$
is the bivariant class which preserves these direct
product decompositions and acts by the already defined
operations $c_i(\mathcal{E}|_{X_r}) \cap -$
on the factors. Observe that in this setting it may happen
that $c_p(\mathcal{E})$ is nonzero for infinitely many $p$.
It follows that the total chern class is an element
$$
c(\mathcal{E}) =
c_0(\mathcal{E}) + c_1(\mathcal{E}) + c_2(\mathcal{E}) + \ldots
\in A^*(X)^\wedge
$$
of the completed bivariant cohomology ring, see
Remark \ref{remark-completion-bivariant}.
In this setting we define the ``rank'' of $\mathcal{E}$
to be the element $r(\mathcal{E}) \in A^0(X)$
as the bivariant operation which sends $(\alpha_r) \in \prod \CH_*(X'_r)$
to $(r\alpha_r) \in \prod \CH_*(X'_r)$.
Note that it is still true that $c_p(\mathcal{E})$ and $r(\mathcal{E})$
are in the center of $A^*(X)$.
\end{remark}

\begin{remark}
\label{remark-top-chern-class}
Let $(S, \delta)$ be as in Situation \ref{situation-setup}. Let $X$
be locally of finite type over $S$. Let $\mathcal{E}$ be a
finite locally free $\mathcal{O}_X$-module. In general
we write $X = \coprod X_r$ as in
Remark \ref{remark-extend-to-finite-locally-free}.
If only a finite number of the $X_r$ are nonempty, then
we can set
$$
c_{top}(\mathcal{E}) = \sum\nolimits_r c_r(\mathcal{E}|_{X_r})
\in A^*(X) = \bigoplus A^*(X_r)
$$
where the equality is Lemma \ref{lemma-disjoint-decomposition-bivariant}.
If infinitely many $X_r$ are nonempty, we will use the same
notation to denote
$$
c_{top}(\mathcal{E}) = \prod c_r(\mathcal{E}|_{X_r})
\in \prod A^r(X_r) \subset A^*(X)^\wedge
$$
see Remark \ref{remark-completion-bivariant} for notation.
\end{remark}











